\SourceLeaf{437}%
\[
\begin{aligned}
d\xi&=\xi'd\lambda\\
d\eta&=\eta'd\lambda\\
d\zeta&=\zeta'd\lambda.
\end{aligned}
\]
By analogy we now call
\[
\frac{d\lambda}{ds}
\]
the measure of curvature of the curved line on the curved surface, and its reciprocal
\[
\frac{ds}{d\lambda}
\]
the radius of curvature. If we denote the latter by $\varrho$, we have
\[
\begin{aligned}
\varrho d\xi&=\xi'ds\\
\varrho d\eta&=\eta'ds\\
\varrho d\zeta&=\zeta'ds.
\end{aligned}
\]

If our line is a shortest line, then $\xi',\eta',\zeta'$ must be proportional to the quantities $X,Y,Z$. But since at the same time
\[
\xi'\xi'+\eta'\eta'+\zeta'\zeta'=XX+YY+ZZ=1,
\]
it follows that
\[
\xi'=\pm X,\qquad\eta'=\pm Y,\qquad\zeta'=\pm Z,
\]
and since furthermore
\[
\begin{aligned}
\xi'X+\eta'Y+\zeta'Z&=\cos\lambda'L\\
&=\pm(XX+YY+ZZ)\\
&=\pm1
\end{aligned}
\]
holds and we always choose the point $\lambda'$ so that
\[
\lambda'L<90^{0}
\]
it follows that for the shortest line
\[
\lambda'L=0
\]
must hold, or $\lambda'$ must coincide with $L$. Thus
